\section{Thesis Structure}
\textbf{Chapter 2 - Background} introduces the background information related 
to the following areas: standard radar algorithms, monopulse angle estimation, 
multi-target tracking, and parallel processing.

\textbf{Chapter 3 - Literature Review} conducts and extensive literature 
review of both foundational and state of the art approaches to performant 
radar systems under high SWaP constraints. The methodologies of works with 
similar motivations to ours are critically evaluated for their effectiveness 
in meeting their aims. Key findings and gaps across works are summarised and 
an explanation for how this thesis fits in with the existing literature is 
provided. 

\textbf{Chapter 4 - Methodology} Explains the methodology by which we compared 
different radar processing pipeline designs for their suitability for use in 
detect and avoid systems, and compare different algorithms for their ability 
to detect small, faraway targets. This methodology will describe the 
algorithms, the implementations of these algorithms, the processes by which 
the implementations are profiled, and the processes by which the algorithms 
are compared for their impact on the overall sensitivity of the system. 

\textbf{Chapter 5 - Results} Presents the results associated with each of the 
experiments outlined in the methodology. The results are illustrated as tables 
and figures, with descriptions of the findings of each of the methodology's 
stages.

\textbf{Chapter 6 - Discussion} The key findings in the results are discussed. 
The results of different experiment stages are evaluated and explained. 
Recommendations are provided. The results are tied in with the problem 
and contributions introduce in this chapter and Chapter 3. The contributions 
that were not achieved are reflected on. 

\textbf{Chapter 7 - Project Plan} A summary of what was achieved so far is provided, and 
a plan for the remaining work is outlined. The plan includes a timeline, stakeholder management plan, and 
a risk management plan.   
